\section{Self-Supervised Objectives}
\label{sec:ssl-objectives}

A self-supervised (SSL) objective trains the encoder to predict a target derived from the graph itself, without human labels. Under missing features it is appealing for a specific reason: the target is computed from information that survives missingness (graph structure, observed neighbors, or observed labels), so the objective can shape the representations of nodes whose own features are absent. By forcing the encoder to solve an auxiliary prediction task, SSL can inject signal that pure classification, trained only on observable nodes, does not reach.

We organize the design space along four axes.

\paragraph{(1) Type (primary axis).} The source of the prediction target, which is our main categorization. \emph{Feature-related} objectives predict statistics of a node's neighbors' features, such as the neighborhood mean or spread, or reconstruct masked feature values; they assume the features themselves carry the signal. \emph{Structural} objectives predict topological quantities computed purely from the graph, such as reachability between nodes, hop distances, landmark distances, or centrality, and thus remain defined even when no features are observed. \emph{Label-related} objectives predict a target built from the observed training labels, for instance the class distribution of a node's neighbors; they are semi-supervised, using labels available at training time but requiring none at inference. \emph{Contrastive} objectives do not predict a fixed target at all: they encode two views of the graph and encourage agreement (or reduced redundancy) between the resulting representations.

\paragraph{(2) Locality.} The receptive field of the target. A \emph{local} target depends only on a node's immediate neighborhood, its adjacent nodes and their features or labels, and is therefore recoverable by a few rounds of message passing. A \emph{global} target depends on a node's position in the graph as a whole, such as its distance to distant landmark nodes or its centrality, and captures long-range structure that local aggregation does not directly expose.

\paragraph{(3) Static vs.\ stochastic.} Whether the target is fixed for the whole of training or refreshed each step. A \emph{static} target is computed once and reused every epoch, so the objective repeatedly presents the same supervision. A \emph{stochastic} target is resampled at each step, for example by drawing fresh random walks, sampling new node pairs, or reselecting landmark nodes, so the objective presents a continually changing supervision signal drawn from the same underlying distribution.

\paragraph{(4) Input-varying vs.\ fixed-input.} What the objective takes as input. A \emph{fixed-input} (per-node) objective maps a single node's representation to a target that is a fixed function of that node; the encoder sees the same input--target pair on every visit, so once it fits each node the objective stops providing new information. An \emph{input-varying} (relational) objective takes a freshly sampled set of nodes and predicts a relation among them, for example whether two sampled nodes are within a few hops or which of two nodes is closer to a landmark, so each step poses a genuinely new problem rather than repeating a memorized one. Note this is distinct from axis (3): a target can be stochastically resampled yet still be per-node (fixed-input), or relational.

Table~\ref{tab:ssl-taxonomy} places representative objectives on these axes, grouped by type.

\paragraph{Target difficulty.} Objectives also differ in how hard their target is to fit. An objective whose target has little variance across nodes, such as the mean of neighbor features that message passing tends to produce anyway, is easy: the loss falls to a small value within a few epochs and then stays flat. An objective whose target has high entropy or depends on a freshly sampled set of nodes each step, such as predicting a label distribution or a relation between sampled nodes, is hard: the loss keeps decreasing over training as the encoder is repeatedly asked new questions. Difficulty therefore follows from the axes above: static, low-variance, feature-space targets tend to be easy, whereas high-entropy or input-varying targets tend to be hard.

\begin{table}[t]
\centering
\caption{Taxonomy of self-supervised objectives, grouped by type. ``Target'' names
the predicted quantity; ``Input'' is per-node (fixed) or relational (varying).}
\label{tab:ssl-taxonomy}
\small
\begin{tabular}{@{}llllll@{}}
\toprule
\textbf{Type} & \textbf{Objective} & \textbf{Target} & \textbf{Locality} & \textbf{Static/Stoch.} & \textbf{Input} \\
\midrule
\multirow{4}{*}{\shortstack[l]{Feature-\\related}}
  & centroid        & neighbor mean            & local  & static     & per-node \\
  & stats           & neighbor std             & local  & static     & per-node \\
  & recon           & masked node features     & local  & stochastic & per-node \\
  & denoising       & denoised features        & local  & stochastic & per-node \\
\cmidrule(l){2-6}
\multirow{9}{*}{Structural}
  & path            & multi-hop reachability   & local  & stochastic & relational \\
  & triplet         & distance ranking         & local  & stochastic & relational \\
  & common-neighbor & pairwise overlap score   & local  & stochastic & relational \\
  & random-walk PE  & return probabilities     & local\,--\,mid & static/stoch. & per-node \\
  & anchor          & distance to landmarks    & global & stochastic & per-node / rel. \\
  & anchorcls       & landmark-distance bucket & global & stochastic & per-node \\
  & shortest-path   & pairwise hop distance    & global & stochastic & relational \\
  & PageRank        & global importance        & global & static     & per-node \\
  & k-core / betw.  & density / role           & global & static     & per-node \\
\cmidrule(l){2-6}
\multirow{3}{*}{\shortstack[l]{Label-\\related}}
  & hist            & neighbor label histogram & local  & static     & per-node \\
  & pseudo-label    & propagated pseudo-labels & local  & stochastic & per-node \\
  & label-smooth    & label smoothness         & local  & static     & per-node \\
\cmidrule(l){2-6}
\multirow{5}{*}{Contrastive}
  & Barlow-Twins      & view decorrelation     & global & stochastic & relational \\
  & InfoNCE / GRACE   & node view agreement    & local  & stochastic & relational \\
  & prototype-InfoNCE & class-prototype agree. & global & stochastic & relational \\
  & DGI / GMI         & node--summary MI       & global\,/\,local & stochastic & relational \\
  & BYOL / AFGRL      & bootstrap agreement    & local  & stochastic & relational \\
\bottomrule
\end{tabular}
\end{table}

In this work we implement and evaluate a representative subset spanning all four types: \emph{feature-related} (centroid, stats, recon); \emph{structural} (path, triplet, anchor, and its classification variant anchorcls); \emph{label-related} (hist); and \emph{contrastive} (Barlow-Twins, InfoNCE, and prototype-InfoNCE). The remaining objectives in Table~\ref{tab:ssl-taxonomy} are included to situate these choices within the broader design space. Formal definitions of the evaluated objectives are given in Appendix~\ref{app:objectives}.
